\documentclass[11pt,letterpaper]{book}
\usepackage[utf8]{inputenc}
\usepackage[T1]{fontenc}
\usepackage[english]{babel}
\usepackage{amsmath,amssymb,amsthm}
\usepackage{graphicx}
\usepackage{float}
\usepackage{hyperref}
\usepackage{geometry}
\usepackage{fancyhdr}
\usepackage{listings}
\usepackage{xcolor}
\usepackage{tikz}
\usepackage{pgfplots}
\usepackage{booktabs}
\usepackage{multirow}
\usepackage{subcaption}
\usepackage{enumitem}
\usepackage{algorithm}
\usepackage{algpseudocode}

% Page geometry
\geometry{margin=1in}

% Colors for code
\definecolor{codegreen}{rgb}{0,0.6,0}
\definecolor{codegray}{rgb}{0.5,0.5,0.5}
\definecolor{codepurple}{rgb}{0.58,0,0.82}
\definecolor{backcolour}{rgb}{0.95,0.95,0.92}

% Code listings setup
\lstdefinestyle{mystyle}{
    backgroundcolor=\color{backcolour},
    commentstyle=\color{codegreen},
    keywordstyle=\color{magenta},
    numberstyle=\tiny\color{codegray},
    stringstyle=\color{codepurple},
    basicstyle=\ttfamily\footnotesize,
    breakatwhitespace=false,
    breaklines=true,
    captionpos=b,
    keepspaces=true,
    numbers=left,
    numbersep=5pt,
    showspaces=false,
    showstringspaces=false,
    showtabs=false,
    tabsize=2
}
\lstset{style=mystyle}

% Headers and footers
\pagestyle{fancy}
\fancyhf{}
\fancyhead[L]{\leftmark}
\fancyhead[R]{\thepage}
\fancyfoot[C]{\textit{Max Dama 2026: State-of-the-Art Quantitative Trading}}

% Custom theorem environments
\newtheorem{theorem}{Theorem}[chapter]
\newtheorem{lemma}[theorem]{Lemma}
\newtheorem{proposition}[theorem]{Proposition}
\newtheorem{corollary}[theorem]{Corollary}
\newtheorem{definition}[theorem]{Definition}
\newtheorem{example}[theorem]{Example}
\newtheorem{remark}[theorem]{Remark}

% Custom commands
\newcommand{\R}{\mathbb{R}}
\newcommand{\N}{\mathbb{N}}
\newcommand{\Z}{\mathbb{Z}}
\newcommand{\Q}{\mathbb{Q}}
\newcommand{\E}{\mathbb{E}}
\newcommand{\Var}{\mathrm{Var}}
\newcommand{\Cov}{\mathrm{Cov}}
\newcommand{\Corr}{\mathrm{Corr}}
\newcommand{\Prob}{\mathbb{P}}
\newcommand{\IC}{\mathrm{IC}}
\newcommand{\IR}{\mathrm{IR}}
\newcommand{\Sharpe}{\mathrm{Sharpe}}
\newcommand{\MaxDD}{\mathrm{MaxDD}}
\newcommand{\OFI}{\mathrm{OFI}}
\newcommand{\VPIN}{\mathrm{VPIN}}
\newcommand{\GLFT}{\mathrm{GLFT}}

\title{Max Dama 2026: \\
Automated Trading at the Frontier \\
\textit{A Journey from Basics to State-of-the-Art Quantitative Strategies}}

\author{Sven Benson \\
\textit{Inspired by Max Dama, Powered by Q23 Quant System}}

\date{2026}

\begin{document}

\maketitle

\tableofcontents

\chapter*{Preface}
\addcontentsline{toc}{chapter}{Preface}

Welcome to \textit{Max Dama 2026: Automated Trading at the Frontier}. This book represents a quantum leap forward from traditional quantitative trading education. While Max Dama pioneered the art of making quant trading accessible and engaging, we now stand at the frontier of 2026, where artificial intelligence, behavioral finance, and market microstructure converge to create strategies that were once the domain of legendary hedge funds.

\textbf{What Makes This Different?}

This is not just another trading book. This is a comprehensive journey that transforms you from a curious beginner into a sophisticated quant practitioner, armed with the same cutting-edge techniques that power the world's most successful systematic traders.

\textbf{Our Approach:}
\begin{itemize}
    \item \textbf{Methodical Progression:} We build from first principles, ensuring no concept is taken for granted
    \item \textbf{State-of-the-Art Techniques:} Every strategy and factor is drawn from current research and proven in live markets
    \item \textbf{Practical Implementation:} All concepts are implemented in production-ready code
    \item \textbf{Risk-Aware Development:} We emphasize robustness, drawdown management, and real-world constraints
\end{itemize}

\textbf{Who This Book Is For:}
\begin{itemize}
    \item Aspiring quantitative traders seeking to understand modern systematic strategies
    \item Finance professionals wanting to bridge the gap between theory and practice
    \item Technology professionals interested in financial applications of machine learning
    \item Students of quantitative finance looking for cutting-edge research applications
\end{itemize}

\textbf{Acknowledgments:}

This work builds upon the foundational work of Max Dama, who first made quantitative trading accessible and fun. We extend our gratitude to the researchers at AQR, WorldQuant, and Two Sigma whose innovations in factor investing and market microstructure inform our approach. Special thanks to the Q23 development team for creating the sophisticated infrastructure that powers these strategies.

\textbf{Ready to Begin?}

Let's embark on this journey together. By the end of this book, you'll not only understand state-of-the-art quantitative trading but also have the tools to implement and extend these techniques yourself.

\chapter{Introduction: The Evolution of Quantitative Trading}

\section{The Quant Revolution}

Quantitative trading has evolved from simple statistical arbitrage in the 1980s to today's sophisticated blend of machine learning, behavioral finance, and market microstructure analysis. What began as academic curiosity has become the dominant force in global markets, managing trillions of dollars in assets.

\section{Our Journey: From Basics to the Frontier}

This book takes you on a comprehensive journey:
\begin{enumerate}
    \item \textbf{Foundations:} Probability, statistics, and financial mathematics
    \item \textbf{Traditional Factors:} Momentum, value, quality, and volatility
    \item \textbf{Market Microstructure:} Order flow, liquidity, and execution quality
    \item \textbf{Behavioral Finance:} Psychological biases and their quantitative exploitation
    \item \textbf{Machine Learning:} Neural networks, reinforcement learning, and feature engineering
    \item \textbf{Advanced Strategies:} Multi-factor models, risk parity, and portfolio optimization
    \item \textbf{Implementation:} Production systems, risk management, and performance analysis
\end{enumerate}

\section{The Q23 Framework}

Throughout this book, we'll work within the Q23 quantitative framework, a state-of-the-art system that implements these concepts in production. The Q23 system features:

\begin{itemize}
    \item \textbf{24 Advanced Strategies:} From simple factor models to neural network-based approaches
    \item \textbf{40+ Factors:} Traditional and novel signals spanning microstructure, behavioral finance, and ML
    \item \textbf{Real-time Analytics:} Live performance monitoring and risk assessment
    \item \textbf{Modular Architecture:} Easily extensible for research and development
\end{itemize}

\begin{figure}[H]
\centering
\begin{tikzpicture}[scale=0.6]
    % Data Layer
    \node[draw, fill=blue!10, rounded corners, minimum height=1cm, minimum width=4cm] at (0,4) {Data Layer};
    \node[draw, fill=blue!20, rounded corners, minimum width=3cm] at (0,3.2) {Market Data};
    \node[draw, fill=blue!20, rounded corners, minimum width=3cm] at (0,2.8) {Order Book};
    \node[draw, fill=blue!20, rounded corners, minimum width=3cm] at (0,2.4) {Trade Flow};

    % Strategy Layer
    \node[draw, fill=green!10, rounded corners, minimum height=2cm, minimum width=6cm] at (0,0) {Strategy Layer};
    \node[draw, fill=green!20, rounded corners, minimum width=2cm] at (-2,-0.5) {GLFT v3};
    \node[draw, fill=green!20, rounded corners, minimum width=2cm] at (0,-0.5) {Neural Alpha};
    \node[draw, fill=green!20, rounded corners, minimum width=2cm] at (2,-0.5) {Benchmarks};
    \node[draw, fill=green!20, rounded corners, minimum width=2cm] at (-2,-1.5) {GLFT v1};
    \node[draw, fill=green!20, rounded corners, minimum width=2cm] at (0,-1.5) {Q23 Composer};
    \node[draw, fill=green!20, rounded corners, minimum width=2cm] at (2,-1.5) {Hybrid Alpha};

    % Factor Library
    \node[draw, fill=purple!10, rounded corners, minimum height=1.5cm, minimum width=4cm] at (-4,1) {Factor Library};
    \node[font=\tiny] at (-4,0.8) {40+ Factors};
    \node[font=\tiny] at (-4,0.4) {• Microstructure};
    \node[font=\tiny] at (-4,0) {• Behavioral};
    \node[font=\tiny] at (-4,-0.4) {• Momentum};
    \node[font=\tiny] at (-4,-0.8) {• Risk/Regime};

    % Analytics Layer
    \node[draw, fill=orange!10, rounded corners, minimum height=1cm, minimum width=4cm] at (4,1) {Analytics Layer};
    \node[font=\tiny] at (4,0.6) {Performance};
    \node[font=\tiny] at (4,0.2) {Risk Metrics};
    \node[font=\tiny] at (4,-0.2) {Factor Exposure};
    \node[font=\tiny] at (4,-0.6) {Attribution};

    % Arrows
    \draw[->, thick] (0,2.2) -- (0,1.8);
    \draw[->, thick] (-3,1) -- (-1.5,0.5);
    \draw[->, thick] (3,1) -- (1.5,0.5);
    \draw[->, thick] (0,-2) -- (0,-3);

    % Execution Layer
    \node[draw, fill=red!10, rounded corners, minimum height=0.8cm, minimum width=3cm] at (0,-4) {Execution Layer};
    \node[font=\tiny] at (0,-4) {Order Generation};

    \node at (0,5.5) {\textbf{Q23 System Architecture}};
\end{tikzpicture}
\caption{The Q23 quantitative trading system integrates data, strategies, factors, and analytics in a modular architecture.}
\label{fig:q23_architecture}
\end{figure}

\chapter{Mathematical Foundations}

\section{Mathematical Utilities: Q23 Math Utils}

The Q23 framework provides robust mathematical utilities for production-grade quantitative computing. These utilities handle edge cases, numerical stability, and cross-platform compatibility.

\subsection{Robust Z-Score Normalization}

The Q23 framework implements robust z-score normalization using median and MAD (Median Absolute Deviation):

\begin{equation}
z_i = \frac{x_i - \median(X)}{\MAD(X) + \epsilon}
\end{equation}

Where $MAD(X) = 1.4826 \times \median(|X - \median(X)|)$.

\textbf{Implementation:}
\begin{lstlisting}[language=Python,caption=Q23 Robust Z-Score]
def robust_zscore_xr(x: "xr.DataArray", dim: str = "asset", eps: float = 1e-12, clip: float = 5.0) -> "xr.DataArray":
    med = x.median(dim=dim, skipna=True)
    mad = np.abs(x - med).median(dim=dim, skipna=True)
    z = (x - med) / (mad * 1.4826 + eps)
    return z.fillna(0).clip(-clip, clip)
\end{lstlisting}

\subsection{Exponentially Weighted Moving Averages}

Q23 implements numerically stable EWMA with bias reduction:

\begin{align}
s_t &= \lambda s_{t-1} + (1-\lambda) x_t \\
w_t &= \lambda w_{t-1} + (1-\lambda) \\
EWMA_t &= \frac{s_t}{w_t + \epsilon}
\end{align}

\subsection{Capped Simplex Projection}

For portfolio optimization, Q23 solves the capped simplex projection:

\begin{align}
\min_{\mathbf{w}} \quad & ||\mathbf{w} - \mathbf{v}||^2 \\
\text{subject to} \quad & \sum w_i = T \\
& l_i \leq w_i \leq u_i \quad \forall i
\end{align}

Using binary search on the Lagrange multiplier.

\subsection{Safe Correlation Computation}

Robust correlation handling zero variance and missing data:

\begin{equation}
\rho = \frac{\Cov(X,Y)}{\sigma_X \sigma_Y + \epsilon}
\end{equation}

With proper NaN handling and numerical guards.

\section{Probability and Statistics}

\subsection{Random Variables and Distributions}

In quantitative trading, we model asset returns as random variables. Let $R_t$ denote the return of an asset at time $t$. We assume:

\begin{align}
R_t &= \mu + \sigma Z_t \\
Z_t &\sim \mathcal{N}(0,1)
\end{align}

Where $\mu$ is the expected return and $\sigma$ is the volatility.

\subsection{Time Series Analysis}

Financial time series exhibit several key properties:

\begin{theorem}[Stationarity]
A time series $\{X_t\}$ is stationary if its statistical properties are invariant under time shifts.
\end{theorem}

\subsection{Statistical Testing}

Hypothesis testing is crucial for factor evaluation:

\begin{definition}[p-value]
The p-value is the probability of observing a test statistic at least as extreme as the one observed, assuming the null hypothesis is true.
\end{definition}

\section{Portfolio Theory}

\subsection{Modern Portfolio Theory}

Harry Markowitz's Nobel Prize-winning framework:

\begin{theorem}[Efficient Frontier]
The efficient frontier represents the set of optimal portfolios that offer the highest expected return for a given risk level.
\end{theorem}

The portfolio return and risk are given by:

\begin{align}
\E[R_p] &= \sum_{i=1}^n w_i \E[R_i] \\
\sigma_p &= \sqrt{\sum_{i=1}^n \sum_{j=1}^n w_i w_j \Cov(R_i, R_j)}
\end{align}

\subsection{Factor Models}

The Fama-French three-factor model:

\begin{equation}
R_i - R_f = \alpha_i + \beta_{i,MKT}(R_M - R_f) + \beta_{i,SMB}SMB + \beta_{i,HML}HML + \epsilon_i
\end{equation}

\chapter{Factor Investing Fundamentals}

\section{The Factor Landscape}

\begin{figure}[H]
\centering
\begin{tikzpicture}[scale=0.7]
    % Main factors
    \node[draw, fill=blue!20, rounded corners, minimum width=2cm] at (0,0) {Market\\($\beta$)};
    \node[draw, fill=green!20, rounded corners, minimum width=2cm] at (-3,-2) {Value\\(B/M)};
    \node[draw, fill=red!20, rounded corners, minimum width=2cm] at (3,-2) {Growth\\(Earnings)};
    \node[draw, fill=yellow!20, rounded corners, minimum width=2cm] at (-3,-4) {Size\\(Cap)};
    \node[draw, fill=purple!20, rounded corners, minimum width=2cm] at (3,-4) {Momentum\\(Returns)};

    % Arrows showing relationships
    \draw[->, thick] (0,0) -- (-2.5,-1.5);
    \draw[->, thick] (0,0) -- (2.5,-1.5);
    \draw[->, thick] (-3,-2) -- (-2.5,-3.5);
    \draw[->, thick] (3,-2) -- (2.5,-3.5);
    \draw[->, thick] (-3,-4) -- (0,-4);
    \draw[->, thick] (3,-4) -- (0,-4);

    % Factor returns
    \node[draw, fill=orange!20, rounded corners] at (0,-6) {Factor Returns};

    % Asset returns
    \draw[->, thick, red] (0,-5.5) -- (0,-7);
    \node[draw, fill=gray!20, rounded corners] at (0,-8) {Asset Returns};

    % Label
    \node at (0,2) {\textbf{Factor Model Structure}};
    \node at (0,1.5) {\scriptsize Fama-French + Momentum};
\end{tikzpicture}
\caption{The factor model decomposes asset returns into systematic factor exposures plus idiosyncratic risk.}
\label{fig:factor_model}
\end{figure}

\section{What Are Factors?}

Factors are systematic sources of risk and return that explain asset price movements. They represent persistent patterns in markets that can be exploited for alpha generation.

\subsection{The Factor Zoo}

Academic research has identified numerous factors:

\begin{itemize}
    \item \textbf{Market:} $\beta$ relative to market portfolio
    \item \textbf{Value:} Book-to-market ratio
    \item \textbf{Size:} Market capitalization
    \item \textbf{Momentum:} Past returns
    \item \textbf{Quality:} Profitability and leverage
    \item \textbf{Volatility:} Idiosyncratic risk
\end{itemize}

\section{Factor Construction}

\subsection{Z-Score Standardization}

Q23 implements multiple z-score variants for different use cases:

\subsubsection{Robust Z-Score (Production Default)}

\begin{equation}
z_{i,t}^{robust} = \frac{x_{i,t} - \median_t(X)}{\MAD_t(X) + \epsilon}
\end{equation}

Where $MAD_t(X) = 1.4826 \times \median_t(|X - \median_t(X)|)$.

\subsubsection{Traditional Z-Score}

\begin{equation}
z_{i,t}^{standard} = \frac{x_{i,t} - \mu_t}{\sigma_t + \epsilon}
\end{equation}

\subsubsection{Cross-Sectional vs Time-Series Z-Scores}

\textbf{Cross-sectional (default):} Standardizes across assets at each time $t$.

\textbf{Time-series:} Standardizes each asset's history (used for regime detection).

\subsection{Winsorization and Robust Statistics}

Q23 implements robust winsorization for outlier handling:

\begin{equation}
x_{winsorized} = \begin{cases}
    P_5 & x < P_5 \\
    x & P_5 \leq x \leq P_{95} \\
    P_{95} & x > P_{95}
\end{cases}
\end{equation}

With adaptive percentile estimation for small universes.

\subsection{Factor Categories and Mathematical Foundations}

\subsubsection{Defensive Factors}

\textbf{Inverse Volatility (inv\_vol):}
\begin{equation}
inv\_vol_{i,t} = \frac{1}{ATR_{i,t} + \epsilon}
\end{equation}

\textbf{Inverse Idiosyncratic Volatility (inv\_idio):}
\begin{equation}
inv\_idio_{i,t} = \frac{1}{\sigma_{i,t}^{residual} + \epsilon}
\end{equation}

\subsubsection{Momentum Factors}

\textbf{Residual Momentum:}
\begin{align}
\beta_{i,t} &= \frac{\Cov(r_{i}, r_{m})}{\Var(r_m)} \\
r_{i,t}^{residual} &= r_{i,t} - \beta_{i,t} r_{m,t} \\
resid\_mom_{i,t} &= SMA(r_{i}^{residual}, window)
\end{align}

\textbf{Multi-Timeframe Momentum:}
\begin{equation}
MTF\_mom_{i,t} = IC_t^{W} \cdot mom_{i,t}^{W} + IC_t^{M} \cdot mom_{i,t}^{M} + IC_t^{Q} \cdot mom_{i,t}^{Q}
\end{equation}

\subsubsection{Market Microstructure Factors}

\textbf{Order Flow Imbalance (OFI):}
\begin{equation}
OFI_{i,t} = \sum_{j=1}^{n} (P_{i,t}^{buy,j} - P_{i,t}^{sell,j})
\end{equation}

\textbf{Volume-Synchronized PIN (VPIN):}
\begin{equation}
VPIN_t = \frac{\sum_{j=1}^{n} |V_{buy,j} - V_{sell,j}|}{\sum_{j=1}^{n} (V_{buy,j} + V_{sell,j}) + \epsilon}
\end{equation}

\subsubsection{Neural Alpha Factors}

\textbf{Attention Momentum:}
\begin{equation}
attention_{i,t} = \sign(r_{i,t}) \cdot (vol\_ratio_{i,t} - 1) \cdot |z_{i,t}^{return}|
\end{equation}

\textbf{Momentum Quality Ratio (Signal-to-Noise):}
\begin{equation}
MQR_{i,t} = \frac{\E[r_{i}]}{\sigma_{i} + \epsilon} \times \sqrt{252}
\end{equation}

\textbf{Efficiency Ratio (Kaufman):}
\begin{equation}
ER_{i,t} = \frac{|P_t - P_{t-window}|}{\sum_{k=1}^{window} |P_{t-k+1} - P_{t-k}| + \epsilon}
\end{equation}

\subsubsection{Ornstein-Uhlenbeck Mean Reversion}

\textbf{OU Parameter Estimation:}
\begin{align}
\log P_{i,t} &= a_i + b_i \log P_{i,t-1} + \epsilon_{i,t} \\
\theta_i &= -\log(|b_i| + \epsilon) \\
\mu_i &= \frac{a_i}{1 - b_i + \epsilon} \\
half\_life_i &= \frac{\log(2)}{\theta_i + \epsilon}
\end{align}

\textbf{OU Z-Score Signal:}
\begin{equation}
z_{i,t}^{OU} = -\frac{\log P_{i,t} - \mu_i}{\sigma_i + \epsilon}
\end{equation}

\chapter{Market Microstructure: The Hidden Market}

\section{Visualizing Market Structure}

\begin{figure}[H]
\centering
\begin{tikzpicture}[scale=0.8]
    % Market participants
    \node[draw, fill=blue!20, rounded corners] at (0,2) {Buyers};
    \node[draw, fill=red!20, rounded corners] at (6,2) {Sellers};

    % Order book
    \draw[thick] (2,0) rectangle (4,4);
    \node at (3,3.5) {\textbf{Order Book}};

    % Bid side
    \node at (2.5,3) {Bid: \$99.95};
    \node at (2.5,2.5) {Bid: \$99.90};
    \node at (2.5,2) {Bid: \$99.85};

    % Ask side
    \node at (3.5,3) {Ask: \$100.05};
    \node at (3.5,2.5) {Ask: \$100.10};
    \node at (3.5,2) {Ask: \$100.15};

    % Spread
    \draw[<->, thick, red] (2.5,2.8) -- (3.5,2.8);
    \node[red] at (3,2.6) {Spread = \$0.10};

    % Market maker
    \node[draw, fill=green!20, rounded corners] at (3,-1) {Market Maker};

    % Arrows showing flow
    \draw[->, thick, blue] (1,2) -- (2.2,2.5);
    \draw[->, thick, red] (5,2) -- (3.8,2.5);

    \draw[->, thick, blue] (3,-0.7) -- (2.5,0.3);
    \draw[->, thick, red] (3,-0.7) -- (3.5,0.3);
\end{tikzpicture}
\caption{Market microstructure shows how orders interact in the limit order book. The spread represents the cost of immediacy.}
\label{fig:market_microstructure}
\end{figure}

\section{Order Flow and Liquidity}

\subsection{Order Flow Imbalance (OFI)}

Order Flow Imbalance measures buying vs selling pressure:

\begin{equation}
OFI_t = \sum_{i=1}^n (P_i^+ - P_i^-)
\end{equation}

Where $P_i^+$ and $P_i^-$ are buy and sell order volumes at price level $i$.

\subsection{Volume-Synchronized Probability of Informed Trading (VPIN)}

Easley et al.'s VPIN model:

\begin{equation}
VPIN = \frac{\sum_{j=1}^n |V_{B,j} - V_{S,j}|}{\sum_{j=1}^n (V_{B,j} + V_{S,j})}
\end{equation}

Where $V_{B,j}$ and $V_{S,j}$ are buy and sell volumes in bucket $j$.

\section{GLFT: The Optimal Execution Framework}

\subsection{The GLFT Formula}

Glosten and Milgrom's Liquidity Trading Formula provides the theoretical foundation for optimal spread:

\begin{equation}
\delta^* = \gamma \sigma^2 \tau + \frac{2}{\gamma} \ln\left(1 + \frac{\gamma}{k}\right)
\end{equation}

Where:
\begin{itemize}
    \item $\gamma$: Risk aversion parameter
    \item $\sigma^2$: Price volatility
    \item $\tau$: Time horizon
    \item $k$: Order arrival rate
\end{itemize}

\subsection{GLFT Strategy Implementation}

Our GLFTv3 strategy implements 24 microstructure factors:

\begin{lstlisting}[language=Python,caption=GLFT Factor Library Structure]
class GLFTv3FactorLibrary(FactorLibrary):
    """Advanced GLFTv3 factor library with cutting-edge microstructure signals."""
    
    # V1 Factors (10): Basic OFI, VPIN, inventory signals
    # V2 Factors (6): Kyle lambda, volume clocks, flow dispersion
    # V3 Factors (8): MTF alignment, regime toxicity, execution quality
    
    def _factor_glft_mtf_ofi_alignment(self) -> "xr.DataArray":
        """Multi-timeframe OFI alignment consensus."""
        # Implementation combines short/med/long OFI signals
\end{lstlisting}

\chapter{Behavioral Finance in Quantitative Strategies}

\section{Cognitive Biases and Market Anomalies}

\subsection{Disposition Effect}

Investors hold losing positions too long and sell winners too early. We quantify this as:

\begin{equation}
Disposition_{i,t} = \frac{UPL_{i,t}}{UPL_{i,t} + RPL_{i,t}}
\end{equation}

Where $UPL$ and $RPL$ are unrealized and realized profits/losses.

\subsection{Anchoring Bias}

Prices tend to anchor to psychological levels. The anchoring factor measures distance from key levels:

\begin{equation}
Anchoring_{i,t} = \min\left( \frac{|P_t - Level_k|}{P_t} \right) \forall k
\end{equation}

\section{Neural Alpha: Behavioral Factors in Practice}

The Q23 Neural Alpha strategy incorporates 15 novel behavioral finance factors:

\begin{lstlisting}[language=Python,caption=Behavioral Finance Factors]
# Behavioral Finance Factors in Neural Alpha
- attention_momentum: Volume-price attention signal
- disposition_alpha: Unrealized P&L vs anchoring price  
- anchoring_bias: Distance from psychological price levels

# Implementation
def _factor_attention_momentum(self) -> "xr.DataArray":
    """Volume-weighted price momentum as attention proxy."""
    vol_weighted_returns = self.returns * self.volume_zscore
    return z_score_cs(vol_weighted_returns.rolling(window=21).mean())
\end{lstlisting}

\chapter{Machine Learning in Factor Discovery}

\section{Feature Engineering}

\subsection{Automatic Factor Generation}

We use automated methods to discover new factors:

\begin{lstlisting}[language=Python,caption=Automated Factor Discovery]
class StrategyGenerator:
    """Generates new trading strategies using evolutionary algorithms."""
    
    def generate_factors(self, base_factors: List[str]) -> List[str]:
        """Create composite factors through mathematical operations."""
        operations = ['+', '-', '*', '/', 'rank', 'zscore']
        # Generate factor combinations systematically
\end{lstlisting}

\subsection{Neural Factor Models}

Deep learning for factor discovery:

\begin{equation}
\mathbf{f}_t = \phi(\mathbf{x}_t; \theta)
\end{equation}

Where $\mathbf{x}_t$ are raw market observables and $\phi$ is a neural network.

\section{Risk Management}

\subsection{Dynamic Volatility Targeting}

Sharpe-optimal volatility scaling:

\begin{equation}
w_{target} = w_{base} \times \frac{\sigma_{target}}{\sigma_{current}}
\end{equation}

\subsection{Drawdown Control}

Risk-off mechanisms activate during drawdowns:

\begin{equation}
risk\_multiplier = \begin{cases}
    1.0 & DD < threshold_1 \\
    0.6 & threshold_1 \leq DD < threshold_2 \\
    0.5 & DD \geq threshold_2
\end{cases}
\end{equation}

\chapter{Advanced Portfolio Construction}

\section{Information Coefficient Weighting}

Dynamic factor weighting based on predictive power:

\begin{equation}
w_f = \frac{IC_f}{\sum |IC_f|}
\end{equation}

Where $IC_f$ is the information coefficient for factor $f$.

\section{Convex Optimization}

\subsection{Modern Portfolio Optimization}

Q23 implements multiple optimization frameworks:

\subsubsection{Maximum Sharpe Ratio Portfolio}

\begin{align}
\max_{\mathbf{w}} \quad & \frac{\mathbf{w}^T \boldsymbol{\mu}}{\sqrt{\mathbf{w}^T \boldsymbol{\Sigma} \mathbf{w}}} \\
\text{subject to} \quad & \mathbf{1}^T \mathbf{w} = 1 \\
& \mathbf{w} \geq \mathbf{0}
\end{align}

\subsubsection{Minimum Variance Portfolio}

\begin{align}
\min_{\mathbf{w}} \quad & \mathbf{w}^T \boldsymbol{\Sigma} \mathbf{w} \\
\text{subject to} \quad & \mathbf{1}^T \mathbf{w} = 1 \\
& \mathbf{w} \geq \mathbf{0}
\end{align}

\subsubsection{Risk Parity}

\begin{equation}
w_i \cdot \frac{\partial \sigma_p}{\partial w_i} = w_j \cdot \frac{\partial \sigma_p}{\partial w_j} \quad \forall i,j
\end{equation}

\subsection{Information Coefficient Weighting}

Dynamic factor weighting based on rolling IC performance:

\begin{equation}
IC_{f,t} = \Corr(f_{t-horizon}, r_{t})
\end{equation}

\begin{equation}
w_{f,t} = \frac{e^{\lambda IC_{f,t}}}{\sum_f e^{\lambda IC_{f,t}}}
\end{equation}

Where $\lambda$ controls the concentration of weights.

\subsection{Composite Score Construction}

Q23's factor aggregation uses IC-weighted combination:

\begin{equation}
score_{i,t} = \sum_f w_{f,t} \cdot z_{i,t}^f
\end{equation}

With softmax tilt for position concentration:

\begin{equation}
w_{i,t}^{tilted} = \frac{e^{\alpha z_{i,t}^{score}}}{\sum_j e^{\alpha z_{j,t}^{score}}}
\end{equation}

\subsection{Volatility Targeting}

Dynamic leverage for volatility targeting:

\begin{equation}
leverage_t = \frac{\sigma_{target}}{\sigma_{current,t}}
\end{equation}

With exponential smoothing of realized volatility:

\begin{equation}
\sigma_{t}^{realized} = \sqrt{\lambda \cdot (\sigma_{t-1}^{realized})^2 + (1-\lambda) \cdot r_{t-1}^2}
\end{equation}

\section{Strategy Ensembling}

Combining multiple strategies:

\begin{equation}
w_{ensemble} = \arg\max_{\mathbf{w}} \IR(\sum w_i s_i)
\end{equation}

Where $s_i$ are individual strategy returns and $\IR$ is the information ratio.

\chapter{Implementation and Production Systems}

\section{The Q23 Architecture}

\subsection{Strategy Registry and Base Classes}

Q23 implements a modular strategy framework with automatic registration:

\begin{lstlisting}[language=Python,caption=Q23 Strategy Registry]
@dataclass
class StrategyConfig:
    name: str
    display_name: str
    version: str
    min_date: str
    exchanges: List[str]
    factors: List[str]
    topn: int = 20
    max_pos: float = 0.05
    target_vol: float = 0.15
    tc_bps: float = 5.0
    # ... additional config

class StrategyBase(ABC):
    @classmethod
    @abstractmethod
    def strategy_id(cls) -> str:
        pass
    
    @property
    @abstractmethod
    def config(self) -> StrategyConfig:
        pass
    
    @abstractmethod
    def run(self, **kwargs) -> StrategyArtifacts:
        pass
\end{lstlisting}

\subsection{Factor Library System}

The core factor computation engine:

\begin{lstlisting}[language=Python,caption=Factor Library Architecture]
class FactorLibrary:
    """Factor computation with configurable windows."""
    
    def __init__(self, market_data, *, params: FactorParams = None):
        self.params = params or FactorParams()
        # Precompute common series
        self.close = self.data["close"].fillna(0.0)
        self.ret = _clip_ret(self.close / self.close.shift(time=1) - 1.0)
        self.beta = _rolling_beta(self.ret, self.mkt, self.params.BETA_WIN)
        self.resid = _residual_returns(self.ret, self.beta, self.mkt)
    
    def _cs_z(self, da: "xr.DataArray") -> "xr.DataArray":
        """Cross-sectional z-score normalization."""
        if self.params.robust_cs_z:
            return _robust_zscore(da, dim=self.asset_dim)
        return safe_zscore(da, dim=self.asset_dim)
    
    def compute(self) -> Tuple["xr.DataArray", Dict[str, "xr.DataArray"]]:
        """Compute all factors with automatic z-normalization."""
        # Implementation with error handling and artifacts
\end{lstlisting}

\subsection{Mathematical Kernel Functions}

Core mathematical operations with numerical stability:

\begin{lstlisting}[language=Python,caption=Mathematical Kernel Functions]
# Rolling statistics with xarray
def _sma(x: "xr.DataArray", win: int) -> "xr.DataArray":
    return x.rolling(time=win, min_periods=win//2).mean()

def _std(x: "xr.DataArray", win: int) -> "xr.DataArray":
    return x.rolling(time=win, min_periods=win//2).std()

def _rolling_beta(returns, mkt, win, eps=1e-12):
    r, m = xr.align(returns, mkt, join="inner")
    cov = _sma(r * m, win) - _sma(r, win) * _sma(m, win)
    var = _sma(m * m, win) - _sma(m, win) ** 2
    return (cov / (var + eps)).fillna(0.0)
\end{lstlisting}

\subsection{Factor Parameter Management}

Dynamic window configuration for regime adaptation:

\begin{lstlisting}[language=Python,caption=Factor Window Configuration]
@dataclass(frozen=True)
class FactorWindowConfig:
    """Window parameters for factor computation."""
    BETA_WIN: int = 162
    MOM_WIN: int = 84
    IDIO_WIN: int = 63
    # ... 30+ parameters
    
    def scale(self, factor: float, min_window: int = 5) -> "FactorWindowConfig":
        """Scale all windows for regime adaptation."""
        return FactorWindowConfig(
            BETA_WIN=max(min_window, int(self.BETA_WIN * factor)),
            MOM_WIN=max(min_window, int(self.MOM_WIN * factor)),
            # ... scaled parameters
        )
\end{lstlisting}

\subsection{Factor Library System}

Modular factor implementation:

\begin{lstlisting}[language=Python,caption=Factor Library Architecture]
class FactorLibrary(ABC):
    """Abstract base class for factor computation."""
    
    def compute(self) -> "xr.DataArray":
        """Compute all requested factors."""
        results = {}
        for factor_name in self.factors:
            method_name = f"_factor_{factor_name}"
            if hasattr(self, method_name):
                results[factor_name] = getattr(self, method_name)()
        return xr.DataArray(results)
\end{lstlisting}

\section{Risk Monitoring}

\subsection{Performance Attribution}

\begin{figure}[H]
\centering
\begin{tikzpicture}[scale=0.8]
    \begin{axis}[
        width=10cm,
        height=6cm,
        title={Strategy Performance Attribution},
        xlabel={Time},
        ylabel={Cumulative Return (\%)},
        legend pos=north west,
        grid=major,
        ymin=-20, ymax=80
    ]

    % GLFT v3
    \addplot[blue, thick] coordinates {
        (0,0) (1,5) (2,12) (3,8) (4,15) (5,22) (6,18) (7,25) (8,30) (9,28) (10,35)
    };
    \addlegendentry{GLFT v3}

    % Neural Alpha
    \addplot[red, thick] coordinates {
        (0,0) (1,3) (2,8) (3,12) (4,9) (5,18) (6,22) (7,19) (8,28) (9,32) (10,38)
    };
    \addlegendentry{Neural Alpha}

    % Benchmark
    \addplot[black, dashed] coordinates {
        (0,0) (1,2) (2,4) (3,6) (4,8) (5,10) (6,12) (7,14) (8,16) (9,18) (10,20)
    };
    \addlegendentry{SPY Benchmark}

    % Drawdown periods
    \fill[red!20] (3,8) rectangle (4,15);
    \fill[red!20] (6,18) rectangle (7,25);
    \node[red] at (3.5,16) {\scriptsize Drawdown};

    \end{axis}
\end{tikzpicture}
\caption{Performance comparison showing GLFT v3 and Neural Alpha strategies versus benchmark. Red areas highlight drawdown periods.}
\label{fig:performance_attribution}
\end{figure}

\subsection{Real-time Analytics}

Live performance tracking:

\begin{equation}
\IR_{realized} = \frac{\E[R_p]}{\sigma_p} \times \sqrt{252}
\end{equation}

\subsection{Stress Testing}

\subsubsection{Monte Carlo Simulation for Tail Risk}

Q23 implements comprehensive Monte Carlo stress testing:

\begin{equation}
VaR_\alpha = -\inf\{x : \Prob(R \leq x) \leq 1-\alpha\}
\end{equation}

With bootstrap resampling and regime-aware scenarios.

\subsubsection{Expected Shortfall (CVaR)}

\begin{equation}
ES_\alpha = -\E[R | R \leq -VaR_\alpha]
\end{equation}

\subsubsection{Maximum Drawdown Distribution}

\begin{equation}
MDD = \max_{t \in [0,T]} \left( \frac{P_t - \max_{s \in [0,t]} P_s}{\max_{s \in [0,t]} P_s} \right)
\end{equation}

\subsubsection{Regime Stress Testing}

Q23 implements regime-aware stress testing:

\begin{enumerate}
    \item Historical regime identification via clustering
    \item Regime-conditioned return simulation
    \item Cross-regime transition modeling
    \item Tail event amplification in stressed regimes
\end{enumerate}

\subsection{Dynamic Risk Management}

\subsubsection{Volatility Regime Detection}

\begin{equation}
regime_t = \begin{cases}
    \text{calm} & \sigma_t < Q_{0.25}(\sigma_{252d}) \\
    \text{normal} & Q_{0.25} \leq \sigma_t \leq Q_{0.75} \\
    \text{volatile} & \sigma_t > Q_{0.75}(\sigma_{252d})
\end{cases}
\end{equation}

\subsubsection{Adaptive Position Sizing}

\begin{equation}
position_{i,t} = f(regime_t) \cdot signal_{i,t} \cdot leverage_t
\end{equation}

Where $f(regime_t)$ implements regime-specific scaling.

\subsubsection{Correlation Breakdown Stress}

Q23 models correlation regime shifts:

\begin{equation}
\rho_{stressed} = \rho_{normal} \cdot (1 - stress\_factor \cdot e^{-\lambda t})
\end{equation}

With exponential decay back to normal correlations.

\subsection{Performance Attribution}

\subsubsection{Factor Contribution Analysis}

\begin{equation}
contribution_{f,t} = w_{f,t} \cdot IC_{f,t} \cdot \sigma_{factor,f}
\end{equation}

\subsubsection{Risk Decomposition}

\begin{equation}
\sigma_p^2 = \sum_f w_f^2 \sigma_f^2 + 2 \sum_{f<g} w_f w_g \sigma_f \sigma_g \rho_{fg} + \sigma_{\epsilon}^2
\end{equation}

\subsubsection{Strategy Alpha/Beta Decomposition}

\begin{align}
\alpha &= \E[R_p - R_b] - \beta \cdot \E[R_m - R_f] \\
\beta &= \frac{\Cov(R_p, R_m)}{\Var(R_m)}
\end{align}

\chapter{Looking Forward: The Future of Quant Trading}

\section{The Evolution of Quantitative Strategies}

\begin{figure}[H]
\centering
\begin{tikzpicture}[scale=0.8]
    % Timeline
    \draw[thick] (0,0) -- (12,0);
    \node at (0,-0.5) {1980s};
    \node at (3,-0.5) {1990s};
    \node at (6,-0.5) {2000s};
    \node at (9,-0.5) {2010s};
    \node at (12,-0.5) {2020s+};

    % Strategy evolution
    \node[draw, fill=blue!20, rounded corners] at (1,2) {Pairs Trading};
    \node[draw, fill=blue!20, rounded corners] at (1,1) {Index Arb};

    \node[draw, fill=green!20, rounded corners] at (4,3) {Factor Models};
    \node[draw, fill=green!20, rounded corners] at (4,2) {Value Investing};
    \node[draw, fill=green!20, rounded corners] at (4,1) {Momentum};

    \node[draw, fill=red!20, rounded corners] at (7,4) {High Frequency};
    \node[draw, fill=red!20, rounded corners] at (7,3) {Market Making};
    \node[draw, fill=red!20, rounded corners] at (7,2) {Statistical Arb};
    \node[draw, fill=red!20, rounded corners] at (7,1) {Order Flow};

    \node[draw, fill=purple!20, rounded corners] at (10,5) {Machine Learning};
    \node[draw, fill=purple!20, rounded corners] at (10,4) {Reinforcement Learning};
    \node[draw, fill=purple!20, rounded corners] at (10,3) {Neural Factors};
    \node[draw, fill=purple!20, rounded corners] at (10,2) {Behavioral Finance};
    \node[draw, fill=purple!20, rounded corners] at (10,1) {Multi-Asset};

    \node[draw, fill=orange!20, rounded corners] at (13,6) {Quantum Computing};
    \node[draw, fill=orange!20, rounded corners] at (13,5) {Generative AI};
    \node[draw, fill=orange!20, rounded corners] at (13,4) {Real-time Adaptation};
    \node[draw, fill=orange!20, rounded corners] at (13,3) {Cross-Market Integration};
    \node[draw, fill=orange!20, rounded corners] at (13,2) {Causal Inference};
    \node[draw, fill=orange!20, rounded corners] at (13,1) {Systemic Risk Control};

    % Arrows showing progression
    \draw[->, thick] (2,1.5) -- (3.5,2.5);
    \draw[->, thick] (5,2) -- (6.5,2.5);
    \draw[->, thick] (8,2.5) -- (9.5,3);
    \draw[->, thick] (11,3) -- (12.5,4);

    \node at (6,6) {\textbf{Evolution of Quantitative Trading Strategies}};
    \node[font=\tiny] at (6,5.5) {From simple arbitrage to AI-driven adaptive systems};
\end{tikzpicture}
\caption{The evolution of quantitative trading from 1980s index arbitrage to 2020s AI-driven systems.}
\label{fig:strategy_evolution}
\end{figure}

\section{Emerging Technologies}

\subsection{Machine Learning Integration}

\subsubsection{Neural Factor Discovery}

Q23 implements automated factor discovery using deep learning:

\begin{equation}
\mathbf{f}_t = \phi(\mathbf{x}_t; \theta)
\end{equation}

Where $\mathbf{x}_t$ includes price, volume, order book, and alternative data.

\subsubsection{Attention Mechanisms for Factor Importance}

\begin{equation}
\alpha_{i,j} = \frac{\exp(e_{i,j})}{\sum_k \exp(e_{i,k})}
\end{equation}

With factor attention weights for dynamic portfolio construction.

\subsubsection{Transformer Architecture for Time Series}

\begin{equation}
\mathbf{h}_t = \Transformer(\mathbf{x}_{t-n:t}, \mathbf{h}_{t-1})
\end{equation}

For sequence modeling of market dynamics.

\subsection{Reinforcement Learning}

\subsubsection{Portfolio Management as RL}

Q23 formulates portfolio optimization as a Markov Decision Process:

\begin{align}
S_t &= \{\mathbf{w}_{t-1}, \mathbf{r}_t, \mathbf{f}_t, \sigma_t\} \\
A_t &= \Delta\mathbf{w}_t \\
R_t &= r_{p,t} - \lambda \cdot |A_t| - \mu \cdot \max(0, \sigma_t - \sigma_{target})
\end{align}

\subsubsection{Deep Q-Learning for Trade Execution}

\begin{equation}
Q(s,a) = r + \gamma \max_{a'} Q(s',a')
\end{equation}

For optimal trade timing and sizing.

\subsubsection{Policy Gradient Methods}

\begin{equation}
\nabla_\theta J(\theta) = \E_{\pi_\theta} [\nabla_\theta \log \pi_\theta(a|s) \cdot A(s,a)]
\end{equation}

For continuous action spaces in portfolio rebalancing.

\subsection{Generative AI and Strategy Discovery}

\subsubsection{AI-Powered Strategy Generation}

Q23 integrates LLMs for strategy ideation:

\begin{lstlisting}[language=Python,caption=Q23 AI Strategy Generation]
class AIStrategyGenerator:
    """Uses large language models for systematic strategy discovery."""
    
    def generate_from_factors(self, factors: List[str]) -> str:
        """Generate strategy logic from factor combinations."""
        prompt = f"Create a trading strategy using: {', '.join(factors)}"
        return self.llm.generate(prompt)
    
    def optimize_parameters(self, strategy_code: str) -> Dict[str, float]:
        """Use evolutionary algorithms to optimize strategy parameters."""
        return self.evolutionary_optimizer.optimize(strategy_code)
    
    def backtest_generated_strategy(self, strategy_spec: Dict) -> Dict:
        """Execute generated strategy with safety bounds."""
        return self.safe_executor.run(strategy_spec)
\end{lstlisting}

\subsubsection{Automated Factor Engineering}

\begin{equation}
\mathbf{f}_{new} = \mathcal{F}(\mathbf{f}_{existing}; \phi)
\end{equation}

Where $\mathcal{F}$ is a learned transformation network.

\subsubsection{Knowledge Graph for Financial Relationships}

\begin{equation}
P(edge | entities) = \sigma(\mathbf{W} [\mathbf{e}_1; \mathbf{e}_2; \mathbf{r}])
\end{equation}

For modeling complex financial relationships.

\subsection{Quantum Computing Applications}

\subsubsection{Quantum Portfolio Optimization}

\begin{equation}
\min_{\mathbf{w}} \mathbf{w}^T \boldsymbol{\Sigma} \mathbf{w} \quad \text{s.t.} \quad \mathbf{1}^T \mathbf{w} = 1
\end{equation}

Solved using quantum annealing for high-dimensional problems.

\subsubsection{Quantum Monte Carlo for Risk Simulation}

\begin{equation}
\E[g(X)] \approx \frac{1}{N} \sum_{i=1}^N g(X_i)
\end{equation}

With quantum speedup for variance reduction.

\subsection{Real-time Adaptation Systems}

\subsubsection{Online Learning}

\begin{equation}
\theta_{t+1} = \theta_t - \eta_t \nabla_\theta \mathcal{L}(y_t, f(\mathbf{x}_t; \theta_t))
\end{equation}

For adapting to changing market conditions.

\subsubsection{Concept Drift Detection}

\begin{equation}
D_{KL}(P_{old} || P_{new}) = \sum P_{old} \log \frac{P_{old}}{P_{new}}
\end{equation}

For detecting structural breaks in factor relationships.

\subsubsection{Meta-Learning for Strategy Transfer}

\begin{equation}
\theta^* = \arg\min_\theta \E_{\mathcal{T}} [\mathcal{L}_\mathcal{T}(\theta)]
\end{equation}

Where $\mathcal{T}$ represents different market regimes.

\section{Responsible Quant Trading}

\subsection{Market Impact Awareness}

Transaction cost modeling:

\begin{equation}
TC = c \times \sqrt{\frac{V}{ADV}} \times \sigma
\end{equation}

Where $V$ is trade volume and $ADV$ is average daily volume.

\subsection{Systemic Risk Considerations}

Portfolio-level risk monitoring:

\begin{equation}
SystemicRisk = \Corr(P, M) \times \frac{\sigma_P}{\sigma_M}
\end{equation}

\chapter{Conclusion: Your Journey Begins}

You now possess the knowledge and tools to navigate the frontier of quantitative trading. The journey from basic concepts to state-of-the-art strategies has equipped you with:

\begin{itemize}
    \item Deep understanding of factor investing and market microstructure
    \item Practical implementation skills in production trading systems
    \item Risk management frameworks for real-world deployment
    \item Foundation for extending these concepts with emerging technologies
\end{itemize}

Remember: the most successful quant traders are those who combine rigorous mathematics with creative problem-solving and disciplined risk management. The Q23 framework provides the infrastructure; your curiosity and persistence provide the rest.

\textbf{What's Next?}

\begin{enumerate}
    \item Implement the strategies in this book using the Q23 codebase
    \item Backtest and validate on historical data
    \item Paper trade in live markets
    \item Extend with your own research and innovations
    \item Contribute back to the quant community
\end{enumerate}

The future of quantitative trading belongs to those who can bridge theory and practice, mathematics and intuition, technology and finance. Welcome to the frontier.

\appendix

\chapter{Mathematical Appendix}

\section{Essential Quantitative Finance Formulas}

\subsection{Risk-Adjusted Performance Measures}

\subsubsection{Sharpe Ratio}
\begin{equation}
\Sharpe = \frac{\E[R_p - R_f]}{\sigma_p}
\end{equation}

\subsubsection{Information Ratio}
\begin{equation}
\IR = \frac{\E[R_p - R_b]}{\sigma_{p-b}}
\end{equation}

\subsubsection{Sortino Ratio}
\begin{equation}
Sortino = \frac{\E[R_p - R_f]}{\sigma_{downside}}
\end{equation}

\subsubsection{Calmar Ratio}
\begin{equation}
Calmar = \frac{\E[R_p]}{\MaxDD}
\end{equation}

\subsection{Volatility Measures}

\subsubsection{Realized Volatility}
\begin{equation}
\sigma_{realized} = \sqrt{\frac{252}{T} \sum_{t=1}^T r_t^2}
\end{equation}

\subsubsection{EWMA Volatility}
\begin{equation}
\sigma_t^2 = \lambda \sigma_{t-1}^2 + (1-\lambda) r_{t-1}^2
\end{equation}

\subsubsection{Parkinson Volatility (High-Low)}
\begin{equation}
\sigma = \frac{1}{4 \ln 2} \left( \ln \frac{H}{L} \right)^2
\end{equation}

\subsection{Factor Model Mathematics}

\subsubsection{Fama-French Three-Factor Model}
\begin{equation}
R_i - R_f = \alpha_i + \beta_{i,MKT}(R_M - R_f) + \beta_{i,SMB}SMB + \beta_{i,HML}HML + \epsilon_i
\end{equation}

\subsubsection{CAPM Beta}
\begin{equation}
\beta_i = \frac{\Cov(R_i, R_m)}{\Var(R_m)}
\end{equation}

\subsubsection{Information Coefficient}
\begin{equation}
IC = \Corr(\mathbf{f}, \mathbf{r}_{forward})
\end{equation}

\subsection{Market Microstructure}

\subsubsection{Effective Spread}
\begin{equation}
Spread_{effective} = 2 \cdot \frac{|P_{trade} - P_{mid}|}{P_{mid}}
\end{equation}

\subsubsection{Price Impact}
\begin{equation}
Impact = \frac{\Delta P}{Volume} \times \frac{ADV}{MarketCap}
\end{equation}

\subsubsection{Liquidity Measure (Amihud)}
\begin{equation}
ILLIQ = \frac{1}{T} \sum_{t=1}^T \frac{|r_t|}{P_t V_t}
\end{equation}

\subsection{Time Series Analysis}

\subsubsection{Autocorrelation Function}
\begin{equation}
\rho_k = \frac{\sum_{t=k+1}^T (x_t - \bar{x})(x_{t-k} - \bar{x})}{\sum_{t=1}^T (x_t - \bar{x})^2}
\end{equation}

\subsubsection{Augmented Dickey-Fuller Test}
\begin{equation}
\Delta y_t = \alpha + \beta t + \gamma y_{t-1} + \sum_{i=1}^p \delta_i \Delta y_{t-i} + \epsilon_t
\end{equation}

\subsubsection{Half-Life of Mean Reversion}
\begin{equation}
half\_life = \frac{\ln(2)}{-\ln(|b|)}
\end{equation}

Where $b$ is the AR(1) coefficient.

\subsection{Machine Learning Metrics}

\subsubsection{Mean Squared Error}
\begin{equation}
MSE = \frac{1}{n} \sum_{i=1}^n (y_i - \hat{y}_i)^2
\end{equation}

\subsubsection{Root Mean Squared Error}
\begin{equation}
RMSE = \sqrt{MSE}
\end{equation}

\subsubsection{Mean Absolute Error}
\begin{equation}
MAE = \frac{1}{n} \sum_{i=1}^n |y_i - \hat{y}_i|
\end{equation}

\subsubsection{R² Score}
\begin{equation}
R^2 = 1 - \frac{\sum (y_i - \hat{y}_i)^2}{\sum (y_i - \bar{y})^2}
\end{equation}

\subsection{Portfolio Optimization}

\subsubsection{Markowitz Optimization}
\begin{align}
\max_{\mathbf{w}} \quad & \mathbf{w}^T \boldsymbol{\mu} \\
\text{s.t.} \quad & \mathbf{w}^T \boldsymbol{\Sigma} \mathbf{w} = \sigma_{target}^2 \\
& \mathbf{1}^T \mathbf{w} = 1
\end{align}

\subsubsection{Black-Litterman Model}
\begin{equation}
\mu_{BL} = \Pi + \tau \boldsymbol{\Sigma} (\tau \boldsymbol{\Sigma} + \boldsymbol{\Omega})^{-1} (Q - \Pi)
\end{equation}

\subsubsection{Risk Parity}
\begin{equation}
w_i \propto \frac{1}{\sigma_i \sum_j w_j \rho_{ij} \sigma_j}
\end{equation}

\subsection{Drawdown Analysis}

\subsubsection{Maximum Drawdown}
\begin{equation}
\MaxDD = \max_{t \in [0,T]} \left( \frac{P_t - \max_{s \in [0,t]} P_s}{\max_{s \in [0,t]} P_s} \right)
\end{equation}

\subsubsection{Average Drawdown}
\begin{equation}
ADD = \frac{1}{N} \sum_{i=1}^N DD_i
\end{equation}

\subsubsection{Drawdown Duration}
\begin{equation}
Duration = \max(\{t_{recovery} - t_{peak} \})
\end{equation}

\subsection{Transaction Cost Modeling}

\subsubsection{Implementation Shortfall}
\begin{equation}
IS = P_{decision} - P_{execution} + fees
\end{equation}

\subsubsection{Market Impact Model (Almgren)}
\begin{equation}
\Delta P = \sign(Q) \cdot \gamma \cdot \sigma \cdot \left( \frac{|Q|}{ADV} \right)^\beta
\end{equation}

\subsubsection{Spread Cost}
\begin{equation}
Cost_{spread} = \frac{Spread}{2} \times \frac{Volume}{ADV}
\end{equation}

\chapter{Code Appendix}

\section{Q23 System Architecture}

\subsection{Core Mathematical Utilities}

\subsubsection{Math Utils Implementation}

\begin{lstlisting}[language=Python,caption=Q23 Math Utils - Robust Statistics]
def robust_zscore_xr(
    x: "xr.DataArray",
    dim: str = "asset",
    eps: float = 1e-12,
    clip: float = 5.0,
) -> "xr.DataArray":
    """Robust z-score using median/MAD for outlier resistance."""
    med = x.median(dim=dim, skipna=True)
    mad = np.abs(x - med).median(dim=dim, skipna=True)
    z = (x - med) / (mad * 1.4826 + eps)
    return z.fillna(0).clip(-clip, clip)

def ewma_vol_1d(returns: Union[np.ndarray, list], lam: float = 0.95) -> np.ndarray:
    """EWMA volatility with proper initialization."""
    r = np.asarray(returns, dtype=float)
    var = 0.0
    out = np.zeros_like(r, dtype=float)
    for i, rv in enumerate(np.nan_to_num(r, nan=0.0, posinf=0.0, neginf=0.0)):
        var = lam * var + (1.0 - lam) * float(rv) * float(rv)
        out[i] = np.sqrt(max(var, 0.0))
    return out

def project_to_capped_simplex(v, T, lo, hi, eps=1e-12, max_iter=60):
    """Binary search for capped simplex projection."""
    v = np.asarray(v, dtype=float)
    lo = np.asarray(lo, dtype=float)
    hi = np.asarray(hi, dtype=float)

    cap = np.maximum(hi - lo, 0.0)
    T_prime = float(T) - float(np.sum(lo))

    if T_prime <= eps:
        return lo.copy()

    a = v - lo
    tau_lo = float(np.min(a - cap))
    tau_hi = float(np.max(a))

    for _ in range(max_iter):
        tau = 0.5 * (tau_lo + tau_hi)
        y = np.clip(a - tau, 0.0, cap)
        if float(np.sum(y)) > T_prime:
            tau_lo = tau
        else:
            tau_hi = tau

    tau = 0.5 * (tau_lo + tau_hi)
    y = np.clip(a - tau, 0.0, cap)
    w = lo + y

    resid = float(T) - float(np.sum(w))
    if abs(resid) > 1e-10:
        free = (w > lo + eps) & (w < hi - eps)
        if np.any(free):
            w[free] += resid / float(np.sum(free))
            w = np.clip(w, lo, hi)

    return w
\end{lstlisting}

\subsection{Strategy Base Class}

\begin{lstlisting}[language=Python,caption=Q23 Strategy Base Class]
@dataclass
class StrategyConfig:
    name: str
    display_name: str
    version: str
    min_date: str
    exchanges: List[str]
    factors: List[str]
    topn: int = 20
    max_pos: float = 0.05
    target_vol: float = 0.15
    tc_bps: float = 5.0
    long_only: bool = True
    long_seats: int = 20
    short_seats: int = 0
    weight_smooth_alpha: float = 0.30
    score_smooth_win: int = 3
    softmax_tilt_alpha: float = 0.85
    risk_off_stretch: float = 0.60
    risk_off_floor: float = 0.50
    dd_win: int = 252

class StrategyBase(ABC):
    @classmethod
    @abstractmethod
    def strategy_id(cls) -> str:
        pass

    @property
    @abstractmethod
    def config(self) -> StrategyConfig:
        pass

    @abstractmethod
    def run(self, **kwargs) -> StrategyArtifacts:
        pass

    def get_output_dir(self) -> str:
        from q23.shared.config import cfg
        return f"{cfg.paths.OUTPUT_ROOT}/{self.strategy_id()}"
\end{lstlisting}

\subsection{Factor Library Core}

\subsubsection{Factor Library Implementation}

\begin{lstlisting}[language=Python,caption=Q23 Factor Library Core]
class FactorLibrary:
    """Factor computation with PM-configurable windows."""

    def __init__(self, market_data, *, params=None, factors=None, asset_dim="asset"):
        _require_xr()
        self.params = params or FactorParams()
        self.factors = list(factors or V4_24_FACTORS)
        self.asset_dim = asset_dim

        if isinstance(market_data, xr.Dataset):
            self.data = market_data
        else:
            self.data = market_data.data

        # Validate required variables
        for v in ("close", "high", "low", "vol"):
            if v not in self.data:
                raise ValueError(f"Dataset missing required variable '{v}'")

        # Precompute common series
        self.close = self.data["close"].fillna(0.0)
        self.high = self.data["high"].fillna(0.0)
        self.low = self.data["low"].fillna(0.0)
        self.vol = self.data["vol"].fillna(0.0)

        self.ret = _clip_ret(self.close / self.close.shift(time=1) - 1.0)
        self.mkt = _market_return(self.ret)
        self.beta = _rolling_beta(self.ret, self.mkt, self.params.BETA_WIN, eps=self.params.eps)
        self.resid = _residual_returns(self.ret, self.beta, self.mkt)

    def _cs_z(self, da: "xr.DataArray") -> "xr.DataArray":
        """Cross-sectional z-score normalization."""
        if self.params.robust_cs_z:
            return _robust_zscore(da, dim=self.asset_dim, eps=self.params.eps, clip=self.params.zclip)
        return safe_zscore(da, dim=self.asset_dim, eps=self.params.eps, clip=self.params.zclip)

    def compute(self) -> Tuple["xr.DataArray", Dict[str, "xr.DataArray"]]:
        """Compute all factors with automatic z-normalization."""
        Fs = []
        names = []

        for name in self.factors:
            fn = getattr(self, f"_factor_{name}", None)
            if fn is None:
                raise ValueError(f"Unknown factor '{name}'")
            raw = fn().transpose("time", self.asset_dim).fillna(0.0)
            z = self._cs_z(raw).fillna(0.0)
            z.name = name
            Fs.append(z)
            names.append(name)

        F = xr.concat(Fs, dim="factor", join="outer").assign_coords(factor=names)
        F = F.transpose("factor", "time", self.asset_dim)
        F.name = "F"
        return F, dict(self._artifacts)
\end{lstlisting}

\subsubsection{Factor Parameter Configuration}

\begin{lstlisting}[language=Python,caption=Q23 Factor Parameters]
@dataclass(frozen=True)
class FactorParams:
    """Complete factor computation parameters."""
    eps: float = 1e-12
    zclip: float = 5.0
    robust_cs_z: bool = True

    # Window configuration
    BETA_WIN: int = 162
    IDIO_WIN: int = 63
    MOM_WIN: int = 84
    ATR_WIN: int = 42

    # Multi-timeframe
    WEEK: int = 5
    MONTH: int = 21
    QUARTER: int = 63
    YEAR: int = 252
    IC_WINDOW: int = 63

    # GLFT microstructure
    OFI_SHORT_WIN: int = 5
    OFI_MED_WIN: int = 21
    VPIN_WIN: int = 50

    # Neural alpha
    EFFICIENCY_WIN: int = 21
    ATTENTION_WIN: int = 10

    @classmethod
    def from_window_config(cls, window_config: FactorWindowConfig) -> "FactorParams":
        """Create FactorParams from a FactorWindowConfig."""
        return cls(**window_config.to_dict())
\end{lstlisting}

\section{Factor Implementations}

\subsection{Core Factor Mathematics}

\subsubsection{Momentum Factors}

\textbf{Residual Momentum (Q23 Production):}
\begin{equation}
resid\_mom_{i,t} = SMA(r_{i,t}^{residual}, MOM\_WIN)
\end{equation}

Where residual returns are beta-adjusted:
\begin{equation}
r_{i,t}^{residual} = r_{i,t} - \beta_{i,t} \cdot r_{m,t}
\end{equation}

\textbf{Multi-Timeframe Momentum:}
\begin{equation}
MTF\_mom_{i,t} = IC_t^W \cdot mom_{i,t}^W + IC_t^M \cdot mom_{i,t}^M + IC_t^Q \cdot mom_{i,t}^Q
\end{equation}

\textbf{Momentum Quality Factors:}
\begin{align}
MQR_{i,t} &= \frac{SMA(r_i, MOM\_WIN)}{\sigma_i + \epsilon} \times \sqrt{252} \\
persistence_{i,t} &= \Corr(r_{i,t}, r_{i,t-1}; PERSISTENCE\_WIN) \\
divergence_{i,t} &= \frac{resid\_mom_i - price\_mom_i}{\sigma_{divergence} + \epsilon}
\end{align}

\subsubsection{Market Microstructure Factors}

\textbf{Order Flow Imbalance (OFI):}
\begin{equation}
OFI_{i,t} = \sum_{j=1}^n (P_{i,t}^{buy,j} - P_{i,t}^{sell,j})
\end{equation}

\textbf{GLFT Optimal Spread:}
\begin{equation}
\delta^* = \gamma \sigma^2 \tau + \frac{2}{\gamma} \ln\left(1 + \frac{\gamma}{k}\right)
\end{equation}

\textbf{Volume-Synchronized PIN (VPIN):}
\begin{equation}
VPIN_t = \frac{\sum_{j=1}^n |V_{buy,j} - V_{sell,j}|}{\sum_{j=1}^n (V_{buy,j} + V_{sell,j}) + \epsilon}
\end{equation}

\subsubsection{Behavioral Finance Factors}

\textbf{Attention Momentum:}
\begin{equation}
attention_{i,t} = \sign(r_{i,t}) \cdot (vol\_ratio_{i,t} - 1) \cdot |z_{i,t}^{return}|
\end{equation}

\textbf{Disposition Effect Alpha:}
\begin{equation}
disposition_{i,t} = -\frac{P_{i,t} - anchor_{i,t}}{anchor_{i,t} + \epsilon} \cdot recovery_{i,t}
\end{equation}

\subsubsection{Risk and Regime Factors}

\textbf{Volatility of Volatility (VoV):}
\begin{equation}
VoV_{i,t} = \frac{1}{\sigma(\sigma_{i,t}^{return}) + \epsilon}
\end{equation}

\textbf{Regime-Adaptive Momentum:}
\begin{equation}
regime\_mom_{i,t} = w_t \cdot mom_{short} + (1-w_t) \cdot mom_{long}
\end{equation}

Where $w_t = sigmoid(vol\_ratio_t - 1)$.

\subsubsection{Ornstein-Uhlenbeck Mean Reversion}

\textbf{OU Parameter Estimation:}
\begin{align}
\log P_{i,t} &= a_i + b_i \log P_{i,t-1} + \epsilon_{i,t} \\
\theta_i &= -\log(|b_i| + \epsilon) \\
half\_life_i &= \frac{\log(2)}{\theta_i + \epsilon} \\
z_{i,t}^{OU} &= -\frac{\log P_{i,t} - \mu_i}{\sigma_i + \epsilon}
\end{align}

\subsection{Implementation Examples}

\subsubsection{Momentum Factor Implementation}

\begin{lstlisting}[language=Python,caption=Q23 Momentum Factor Implementation]
def _factor_resid_mom(self) -> "xr.DataArray":
    """Residual momentum - beta-adjusted price momentum."""
    return _sma(self.resid, self.params.MOM_WIN).fillna(0.0)

def _factor_momentum_quality_ratio(self) -> "xr.DataArray":
    """Rolling Sharpe-like momentum quality (annualized)."""
    ret_mean = _sma(self.ret, self.params.MOM_WIN)
    ret_std = _std(self.ret, self.params.MOM_WIN)
    mqr = ret_mean / (ret_std + self.params.eps)
    return (mqr * np.sqrt(252)).fillna(0.0).clip(-5, 5)
\end{lstlisting}

\subsubsection{Market Microstructure Implementation}

\begin{lstlisting}[language=Python,caption=GLFT Factor Implementation]
def _compute_ofi(self, window: int) -> "xr.DataArray":
    """Order Flow Imbalance calculation."""
    signed_vol = self._compute_signed_volume()
    vol = self.vol
    signed_sum = signed_vol.rolling(time=window, min_periods=window//2).sum()
    vol_sum = vol.rolling(time=window, min_periods=window//2).sum()
    ofi = signed_sum / (vol_sum + self.eps)
    return ofi.clip(-1, 1)

def _factor_glft_mtf_ofi_alignment(self) -> "xr.DataArray":
    """Multi-timeframe OFI consensus."""
    ofi_s = self._compute_ofi(self.mtf_ofi_short)
    ofi_m = self._compute_ofi(self.mtf_ofi_med)
    ofi_l = self._compute_ofi(self.mtf_ofi_long)
    
    alignment = sign_s + sign_m + sign_l
    magnitude = (np.abs(ofi_s) + np.abs(ofi_m) + np.abs(ofi_l)) / 3
    aligned_signal = alignment * magnitude
    
    return z_score_cs(aligned_signal, self.eps)
\end{lstlisting}

\subsubsection{Neural Alpha Factor Implementation}

\begin{lstlisting}[language=Python,caption=Neural Alpha Factor Implementation]
def _factor_attention_momentum(self) -> "xr.DataArray":
    """Volume-price attention signal."""
    vol_avg = _sma(self.vol, self.params.ADV_WIN_LONG)
    vol_ratio = self.vol / (vol_avg + self.params.eps)
    vol_ratio = vol_ratio.fillna(1.0).clip(0.1, 10.0)
    
    ret_std = _std(self.ret, self.params.MOM_SHORT)
    ret_z = (self.ret / (ret_std + self.params.eps)).fillna(0.0).clip(-3, 3)
    
    attention = np.sign(self.ret) * (vol_ratio - 1.0) * np.abs(ret_z)
    attention_win = getattr(self.params, "ATTENTION_WIN", 10)
    return _sma(attention, int(attention_win)).fillna(0.0)

def _factor_efficiency_ratio(self) -> "xr.DataArray":
    """Kaufman Efficiency Ratio with momentum direction."""
    eff_win = int(getattr(self.params, "EFFICIENCY_WIN", 21))
    
    direction = np.abs(self.close - self.close.shift(time=eff_win))
    abs_changes = np.abs(self.close - self.close.shift(time=1))
    volatility = abs_changes.rolling(time=eff_win, min_periods=eff_win).sum()
    
    er = direction / (volatility + self.params.eps)
    mom_dir = np.sign(self.close - self.close.shift(time=eff_win))
    return (er * mom_dir).fillna(0.0).clip(-1, 1)
\end{lstlisting}

\bibliography{references}
\bibliographystyle{plain}

\end{document}